\scanpage{6}
\textcircled{1.3.1.} Plus \unsure{gén.}, une descente de $\unsure{\mathscr X|\mathcal X|\mathfrak X}$ au corps de base $k'$ revient à un \inserted{$\Gamma'$-}scindage
\diagram{(9)\quad\begin{tikzcd}[ampersand replacement=\&,column sep=small,row sep=small]
\unsure{\mathcal G|\mathscr G}\arrow[r]\&\Gamma\\
\&\Gamma_1\arrow[u,hook]\arrow[ul,dashed]
\end{tikzcd}}

i.e. à un scindage de
\diagram{\begin{tikzcd}[ampersand replacement=\&,column sep=small]
1\arrow[r,dash]\&\unsure{\mathfrak G|G|\mathscr G}\arrow[r,dash]\&\deleted{\unsure{\Gamma}}\mathcal G_1\arrow[r,dash]\&\Gamma_1
\end{tikzcd}}

\unsure{donc|d’où} une écriture correspondante
\[
\mathcal G_1^{\deleted{\unsure{\prime}}}=\Gamma_1^{\deleted{\unsure{\prime}}}\mathbin{\underset{1/2}{\cdot}}\mathfrak G. \tag*{(10)}
\]

\textcircled{2.1} Soit $\widetilde{\mathscr X}$ un rev. universel \inserted{(profini)} de $\mathscr X$, \deleted{\unsure{Si} (plus gén un rev. profini galoisien} \deleted{\unsure{a}} alors une extension \unsure{canonique}, \unsure{où|on}
\[
\unsure{\mathcal E|\mathscr E}=\mathrm{Aut}_k(\widetilde{\mathscr X}\deleted{\unsure{/}}\to\mathscr X)\tag*{(11)}
\]
\margin{\uline{NB} $\mathcal G\to\mathrm{Autext}(\pi_1)$ est au sens des aut. ext. du groupe \uline{profini} $\pi_1$ (i.e. pour $g\in\unsure{\mathcal G|\mathcal E|G}$, $\mathrm{int}(g)\vert\pi_1$ est \uline{continu})}
\diagram{(12)\quad\begin{tikzcd}[ampersand replacement=\&,column sep=small]
1\arrow[r]\&\pi_1\arrow[r]\&\mathcal E\arrow[r]\&\mathcal G\arrow[r,dash]\&1
\end{tikzcd}}

où \unsure{$1$}
\[
\pi_1=\mathrm{Aut}(\widetilde{\mathscr X}/\mathscr X)\tag*{(13)}
\]
est un groupe fondamental de $\mathscr X$ (\unsure{gpe|groupe} profini)

\deleted{Ceci \unsure{reste|vaut} valable \unsure{sans|aussi}}

\textcircled{2.2} Considérons un rev profini galoisien $\mathscr X'$ de $\mathscr X$, correspondant à un ss-gpe fermé invariant
\[
\begin{aligned}
&N\subset\pi_1\qquad \pi_1/N=\deleted{\unsure{H|G}}\deleted{G}H\unsure{\simeq|\leftarrow|=}\mathrm{Aut}(\mathscr X'/\mathscr X)\\
&\mathscr X'=\widetilde{\mathscr X}/N
\end{aligned}\tag*{(12)}
\]
% archive sheet number: 5
